\documentclass[12pt,a4paper]{article}

% Paquetes necesarios
\usepackage[utf8]{inputenc}
\usepackage[T1]{fontenc}
\usepackage[spanish]{babel}
\usepackage{amsmath}
\usepackage{amssymb}
\usepackage{amsthm}
\usepackage{mathtools}
\usepackage{geometry}
\usepackage{hyperref}
\usepackage{bookmark}
\usepackage{graphicx}
\usepackage{booktabs}
\usepackage{array}
\usepackage{enumitem}
\usepackage{fancyhdr}
\usepackage{titlesec}
\usepackage{textcomp}
\usepackage{tikz}
\usetikzlibrary{arrows.meta, positioning, calc, decorations.pathreplacing, shapes.geometric}
\usepackage{pgfplots}
\pgfplotsset{compat=1.18}
\usepgfplotslibrary{fillbetween}
\usepackage{caption}

% Configuración de página
\geometry{margin=2.5cm}
\setlength{\headheight}{15pt}
\pagestyle{fancy}
\fancyhf{}
\fancyhead[L]{Fundamentos de Valoración de Activos}
\fancyhead[R]{Jesus Villota Miranda - CEMFI}
\fancyfoot[C]{\thepage}

% Configuración de hipervínculos
\hypersetup{
    colorlinks=true,
    linkcolor=blue,
    urlcolor=blue,
    citecolor=blue
}

% Definición de entornos matemáticos
\theoremstyle{definition}
\newtheorem{definicion}{Definición}[section]
\newtheorem{teorema}{Teorema}[section]
\newtheorem{proposicion}{Proposición}[section]
\newtheorem{ejemplo}{Ejemplo}[section]
\newtheorem{corolario}{Corolario}[section]
\newtheorem{ejercicio}{Ejercicio}[section]
\newtheorem{solucion}{Solución}[section]

% Comandos personalizados
\newcommand{\R}{\mathbb{R}}
\newcommand{\N}{\mathbb{N}}
\newcommand{\Z}{\mathbb{Z}}
\newcommand{\C}{\mathbb{C}}
\newcommand{\E}{\mathbb{E}}
\newcommand{\V}{\mathbb{V}}
\newcommand{\Cov}{\mathbb{C}\text{ov}}

% Estilos TikZ para visualizaciones financieras
\tikzset{
  timeline/.style={-Stealth, thick, black},
  cashflow/.style={-Stealth, very thick, blue!70},
  discount/.style={-Stealth, dashed, red!70},
  simple/.style={fill=blue!30, draw=blue!70},
  compound/.style={fill=green!30, draw=green!70},
  timepoint/.style={circle, fill, inner sep=2pt},
  payment/.style={-Stealth, thick, orange!80},
  growth/.style={-Stealth, thick, purple!80}
}

% Configuración de títulos
\titleformat{\section}{\Large\bfseries}{\thesection}{1em}{}
\titleformat{\subsection}{\large\bfseries}{\thesubsection}{1em}{}
\titleformat{\subsubsection}{\normalsize\bfseries}{\thesubsubsection}{1em}{}



% Título del documento
\title{\textbf{Fundamentos de Valoración de Activos}\\
\large Parte VII: Modelos de Valoración de Activos con Riesgo}
\author{Jesus Villota Miranda\\
CEMFI}
\date{\today}

\begin{document}

\maketitle

\tableofcontents
\newpage

\section{Introducción}

Este documento presenta un tratamiento riguroso y autocontenido de los modelos de valoración de activos con riesgo, con especial énfasis en el Modelo de Valoración de Activos de Capital (CAPM, por sus siglas en inglés). El CAPM constituye una pieza central de las finanzas modernas y proporciona un marco teórico fundamental para entender cómo se determinan los precios de los activos financieros en equilibrio de mercado.

El documento se estructura de la siguiente manera: primero, se revisan los conceptos fundamentales del modelo de Markowitz y el teorema de separación de fondos, que constituyen la base teórica del CAPM. Posteriormente, se desarrolla el modelo CAPM, incluyendo sus supuestos, la línea del mercado de capitales (CML) y la línea del mercado de títulos (SML). Se analiza en detalle el concepto de beta como medida del riesgo sistemático y su cálculo para activos individuales y carteras. Finalmente, se presentan extensiones y modelos alternativos de valoración, incluyendo estudios de eventos, el modelo de tres factores de Fama-French, el modelo de cuatro factores de Carhart y la Teoría de Valoración por Arbitraje (APT).

\section{Repaso: Modelo de Markowitz y Teorema de Separación de Fondos}

\subsection{El Modelo Media-Varianza}

El modelo de Markowitz, también conocido como modelo media-varianza, constituye un modelo clásico en la gestión de activos que permite a los inversores construir carteras óptimas considerando el trade-off entre rentabilidad esperada y riesgo.

\begin{definicion}[Modelo de Markowitz]
El \textbf{modelo de Markowitz} o \textbf{modelo media-varianza} es un modelo de optimización de carteras que permite a los inversores seleccionar carteras que maximizan su utilidad esperada, considerando únicamente la rentabilidad esperada y la varianza (o desviación típica) de los rendimientos como medidas de rentabilidad y riesgo, respectivamente.
\end{definicion}

Los conceptos clave del modelo de Markowitz incluyen:

\begin{itemize}
\item \textbf{Cartera mínima varianza}: La cartera que minimiza la varianza para un nivel dado de rentabilidad esperada.
\item \textbf{Cartera eficiente}: Una cartera que ofrece la máxima rentabilidad esperada para un nivel dado de riesgo, o el mínimo riesgo para un nivel dado de rentabilidad esperada.
\item \textbf{Frontera eficiente}: El conjunto de todas las carteras eficientes, representado gráficamente como una curva en el espacio rentabilidad-riesgo.
\end{itemize}

\subsection{La Línea de Asignación de Capital (CAL)}

Cuando existe la posibilidad de invertir en un activo libre de riesgo además de activos con riesgo, el conjunto de oportunidades de inversión se expande significativamente.

\begin{definicion}[Línea de Asignación de Capital (CAL)]
La \textbf{línea de asignación de capital} (CAL, Capital Allocation Line) representa todas las carteras que puede formar el inversor combinando una cartera con riesgo (cartera C) y el activo libre de riesgo. La CAL es una línea recta que conecta el activo libre de riesgo con la cartera de riesgo óptima en el espacio rentabilidad-riesgo.
\end{definicion}

\subsection{La Cartera Tangente}

\begin{definicion}[Cartera Tangente]
La \textbf{cartera tangente} es la cartera eficiente que elegirá un inversor si existe la posibilidad de invertir en un activo libre de riesgo. Esta cartera corresponde al punto de tangencia entre la frontera eficiente de activos con riesgo y la línea recta que parte del activo libre de riesgo.
\end{definicion}

\subsection{Teorema de Separación de Fondos}

El teorema de separación de fondos establece un resultado fundamental sobre la estructura óptima de las carteras de los inversores.

\begin{teorema}[Teorema de Separación de Fondos]
Cualquier inversor maximizará su utilidad esperada, con independencia de su comportamiento frente al riesgo, repartiendo su presupuesto únicamente entre el activo libre de riesgo ($R_f$) y la cartera tangente. La diferencia entre inversores radica únicamente en la proporción de su cartera que invierten en la cartera tangente y en el activo libre de riesgo, según sus preferencias de riesgo.
\end{teorema}

\begin{proof}
Dado que todos los inversores tienen acceso a la misma información y siguen el modelo media-varianza, todos identifican la misma cartera tangente como la cartera óptima de activos con riesgo. La elección de la combinación específica entre la cartera tangente y el activo libre de riesgo depende únicamente de las preferencias de riesgo del inversor, representadas por sus curvas de indiferencia. El punto óptimo para cada inversor es aquel en el que su curva de indiferencia es tangente a la CAL.
\end{proof}

La solución del problema de optimización de cartera se puede dividir en dos etapas:

\begin{enumerate}
\item \textbf{Determinar la composición de la cartera tangente T}: Esta etapa es común para todos los inversores y depende únicamente de las características de los activos con riesgo disponibles en el mercado.
\item \textbf{Decidir la proporción a invertir en T y qué parte prestar o pedir prestada a $R_f$}: Esta etapa depende de las preferencias individuales de riesgo de cada inversor.
\end{enumerate}

\section{Ejercicio: Análisis de Cartera con AIR FRANCE y DANONE}

Este ejercicio permite aplicar los conceptos del modelo de Markowitz y el teorema de separación de fondos mediante un ejemplo práctico con dos acciones del CAC40.

\begin{ejemplo}[Análisis Completo de Cartera: AIR FRANCE y DANONE]
Supongamos que un inversor desea formar una cartera invirtiendo únicamente en dos acciones (AIR FRANCE y DANONE) que cotizan en el CAC40. Los datos para estas dos acciones son los siguientes:

\begin{center}
\begin{tabular}{lcc}
\toprule
Activo & $\E[r_i]$ & $\sigma_i$ \\
\midrule
AIR FRANCE & 11\% & 15\% \\
DANONE & 7\% & 12\% \\
Coeficiente de correlación & $\rho = 0.15$ & \\
\bottomrule
\end{tabular}
\end{center}

\textbf{a) Calcule la rentabilidad esperada, varianza y desviación típica de una cartera equiponderada.}

\textbf{Solución:}

Para una cartera equiponderada: $w_1 = w_2 = 0.5$.

Rentabilidad esperada:
$$\E[r_p] = 0.5 \times 0.11 + 0.5 \times 0.07 = 0.09 = 9\%$$

Varianza:
\begin{align}
\sigma_p^2 &= w_1^2 \sigma_1^2 + w_2^2 \sigma_2^2 + 2w_1 w_2 \rho \sigma_1 \sigma_2 \\
&= (0.5)^2 \times (0.15)^2 + (0.5)^2 \times (0.12)^2 + 2 \times 0.5 \times 0.5 \times 0.15 \times 0.15 \times 0.12 \\
&= 0.005625 + 0.0036 + 0.00135 \\
&= 0.010575
\end{align}

Desviación típica:
$$\sigma_p = \sqrt{0.010575} = 0.1029 = 10.29\%$$

\textbf{b) Calcule la rentabilidad esperada y riesgo para diferentes composiciones de cartera.}

Para diferentes pesos, aplicamos las fórmulas del modelo de Markowitz. A continuación se muestran algunos ejemplos de combinaciones:

\begin{itemize}
\item $w_1 = 1.0, w_2 = 0.0$: $\E[r_p] = 11\%$, $\sigma_p = 15\%$
\item $w_1 = 0.9, w_2 = 0.1$: $\E[r_p] = 10.6\%$, $\sigma_p \approx 13.8\%$
\item $w_1 = 0.75, w_2 = 0.25$: $\E[r_p] = 10\%$, $\sigma_p \approx 12.5\%$
\item $w_1 = 0.6, w_2 = 0.4$: $\E[r_p] = 9.4\%$, $\sigma_p \approx 11.2\%$
\item $w_1 = 0.5, w_2 = 0.5$: $\E[r_p] = 9\%$, $\sigma_p = 10.29\%$
\item $w_1 = 0.4, w_2 = 0.6$: $\E[r_p] = 8.6\%$, $\sigma_p \approx 9.8\%$
\item $w_1 = 0.25, w_2 = 0.75$: $\E[r_p] = 8\%$, $\sigma_p \approx 9.5\%$
\item $w_1 = 0.1, w_2 = 0.9$: $\E[r_p] = 7.4\%$, $\sigma_p \approx 9.8\%$
\item $w_1 = 0.0, w_2 = 1.0$: $\E[r_p] = 7\%$, $\sigma_p = 12\%$
\end{itemize}

\textbf{c) Represente gráficamente los resultados e identifique los conceptos clave.}

En el espacio rentabilidad-riesgo $(\sigma_p, \E[r_p])$:

\begin{itemize}
\item \textbf{Conjunto de oportunidades de inversión:} Todos los puntos que representan las diferentes combinaciones de AIR FRANCE y DANONE forman una curva (hipérbola) que conecta los dos activos individuales.

\item \textbf{Carteras dominadas:} Son aquellas carteras que, para un mismo nivel de riesgo, ofrecen menor rentabilidad esperada, o para un mismo nivel de rentabilidad esperada, tienen mayor riesgo. En este caso, las carteras en la parte inferior de la curva están dominadas.

\item \textbf{Cartera de mínima varianza:} Es el punto de la curva con menor riesgo. Se puede calcular usando la fórmula de la cartera de mínima varianza para dos activos.

\item \textbf{Cartera de máxima rentabilidad:} Es el punto con mayor rentabilidad esperada, que corresponde a invertir 100\% en AIR FRANCE ($w_1 = 1$).

\item \textbf{Frontera eficiente:} Es la parte superior de la curva desde la cartera de mínima varianza hasta la cartera de máxima rentabilidad. Solo las carteras en esta frontera son eficientes.
\end{itemize}

\textbf{d) Elección de cartera sin activo libre de riesgo}

Si el inversor es averso al riesgo y no puede acceder a un activo libre de riesgo, elegirá una cartera de la frontera eficiente. La elección específica dependerá de su nivel de aversión al riesgo: cuanto más averso al riesgo, más cerca estará de la cartera de mínima varianza.

Si otro inversor es menos averso al riesgo, elegirá una cartera con mayor rentabilidad esperada y mayor riesgo, más cerca de la cartera de máxima rentabilidad (100\% AIR FRANCE). Ambos elegirán carteras de la frontera eficiente, pero en diferentes puntos según sus preferencias.

\textbf{e) Introducción del activo libre de riesgo}

Si los inversores pueden invertir en Letras del Tesoro con $r_f = 3.5\%$, entonces pueden combinar cualquier cartera con riesgo de la frontera eficiente con el activo libre de riesgo, formando líneas CAL (Capital Allocation Line). Según el teorema de separación de fondos, todos los inversores elegirán la misma cartera con riesgo (la cartera tangente), que es la que maximiza el Ratio de Sharpe. La diferencia entre inversores estará solo en la proporción invertida en la cartera tangente versus el activo libre de riesgo.

\textbf{f) Cálculo de la cartera óptima con $r_f = 3.5\%$}

La cartera tangente es aquella que maximiza el Ratio de Sharpe:
$$\text{Sharpe Ratio} = \frac{\E[r_p] - r_f}{\sigma_p}$$

Para encontrar los pesos óptimos, se debe resolver el problema de optimización que maximiza este ratio sujeto a las restricciones de pesos. La solución requiere resolver un sistema de ecuaciones que iguala las pendientes de las curvas de indiferencia con la pendiente de la CAL en el punto de tangencia.

La cartera tangente tendrá una composición específica de AIR FRANCE y DANONE que maximiza el Ratio de Sharpe, y todos los inversores combinarán esta cartera con el activo libre de riesgo según sus preferencias de riesgo.
\end{ejemplo}

\section{El Modelo de Valoración de Activos de Capital (CAPM)}

\subsection{Supuestos y Origen del CAPM}

El Modelo de Valoración de Activos de Capital (CAPM) es una pieza central de las finanzas modernas, aunque fue desarrollado hace casi medio siglo. El CAPM fue desarrollado por William Sharpe en 1962, trabajo por el cual fue premiado con el Premio Nobel de Ciencias Económicas.

\begin{definicion}[CAPM]
El \textbf{Modelo de Valoración de Activos de Capital} (CAPM, Capital Asset Pricing Model) es un modelo de equilibrio que determina la rentabilidad esperada de un activo en función de su riesgo sistemático, medido por su beta. El modelo parte del modelo media-varianza de Markowitz y asume que el mercado de capitales está en equilibrio, es decir, que la oferta iguala a la demanda.
\end{definicion}

El CAPM es un modelo teórico que requiere de una serie de supuestos para poder desarrollarlo. Estos supuestos son necesarios para simplificar la realidad y permitir la derivación de resultados analíticos precisos.

\begin{teorema}[Supuestos del CAPM]
El modelo CAPM se basa en los siguientes supuestos:

\begin{enumerate}
\item \textbf{Modelo estático}: Los agentes solo miran al próximo periodo (1 trimestre, 1 año, etc.). No se considera un marco temporal multiperiodo.

\item \textbf{Mercado perfectamente competitivo}: Existe una gran cantidad de inversores, cada uno con una función de utilidad y una dotación de riqueza inicial. Los inversores son precio-aceptantes, es decir, no pueden influir en los precios de mercado mediante sus decisiones individuales.

\item \textbf{Oferta exógena y divisibilidad perfecta}: La oferta de los activos financieros con riesgo está dada exógenamente, y estos son perfectamente divisibles. Los inversores pueden comprar fracciones de cualquier activo.

\item \textbf{Tipo de interés único}: El tipo de interés al que se remuneran los fondos es igual al que se paga por disponer de capitales ajenos. No existe diferencia entre el tipo de interés de préstamo y de depósito.

\item \textbf{Ausencia de fricciones}: No existen costes de transacción ni impuestos. Todas las transacciones se realizan sin costes adicionales.

\item \textbf{Optimizadores Markowitz}: Todos los inversores son optimizadores en sentido Markowitz, es decir, solo les interesa la media y la varianza de los rendimientos. Sus funciones de utilidad dependen únicamente de estos dos momentos de la distribución de rendimientos.

\item \textbf{Expectativas homogéneas}: Todo inversor posee igual información e igual datos, y por tanto, sus expectativas de rentabilidad y riesgo para cada activo son idénticas. Todos los inversores tienen las mismas estimaciones de los parámetros de los activos.
\end{enumerate}
\end{teorema}

Estos supuestos, aunque restrictivos, permiten derivar resultados analíticos precisos y proporcionan un marco de referencia fundamental para entender la formación de precios en los mercados financieros.

\section{La Línea del Mercado de Capitales (CML)}

\subsection{Definición y Características}

Dado que todos los inversores poseen igual información y siguen el modelo Media-Varianza, todos los inversores mantienen como cartera con riesgo la cartera tangente (T). La diferencia entre inversores está en la proporción de su cartera que invierten en esa cartera tangente y en el activo libre de riesgo.

\begin{definicion}[Línea del Mercado de Capitales (CML)]
La \textbf{línea del mercado de capitales} (CML, Capital Market Line) es la línea recta que conecta el activo libre de riesgo con la cartera tangente en el espacio rentabilidad-riesgo. La CML representa todas las combinaciones posibles entre el activo libre de riesgo y la cartera tangente.
\end{definicion}

\subsection{La Cartera de Mercado}

Un resultado fundamental del CAPM es que la cartera tangente coincide con la cartera de mercado.

\begin{proposicion}[Equivalencia entre Cartera Tangente y Cartera de Mercado]
Dados los supuestos del CAPM, la composición de la cartera de mercado coincide con la de la cartera tangente.
\end{proposicion}

\begin{proof}
Dado que todos los agentes mantienen como activo con riesgo la cartera tangente (obtenida del modelo media-varianza) y que el CAPM es un modelo en equilibrio (exceso de oferta igual a cero), entonces, necesariamente, el peso de un activo en la cartera tangente será igual al peso de ese activo en la cartera de mercado. Si agregamos todas las carteras con riesgo de todos los agentes de la economía, obtenemos la cartera de mercado (Cartera ``M'').
\end{proof}

\begin{definicion}[Cartera de Mercado]
La \textbf{cartera de mercado} se define como aquella cartera compuesta por todos los activos con riesgo de la economía. Si agregamos todas las carteras con riesgo de todos los agentes de la economía, esta es la cartera de mercado (Cartera ``M'').
\end{definicion}

La proporción de un activo \(j\) en la cartera de mercado será igual al valor total de ese activo \(j\) en la economía dividido por el valor total de todos los activos con riesgo de la economía:

\begin{equation}
w_j = \frac{\text{Valor del activo } j}{\text{Valor de toda la cartera de mercado}} = \frac{n_j P_j^*}{\sum_{i=1}^{N} n_i P_i^*}
\end{equation}

donde $P^*$ representa los precios de equilibrio, tales que Oferta = Demanda.

\begin{ejemplo}[Equivalencia entre Cartera Tangente y Cartera de Mercado]
Suponga que el peso de las acciones de TELEFÓNICA en la cartera tangente es del 23\%. Esto significa que todos los agentes mantienen un 23\% de la riqueza invertida en activos con riesgo en acciones de TELEFÓNICA. Por tanto, si la cartera de mercado es la agregación de todas las carteras de los agentes de la economía, en esta cartera de mercado las acciones de TELEFÓNICA representarán también el 23\%. E igual ocurre con todos los activos.
\end{ejemplo}

\section{La Línea del Mercado de Títulos (SML)}

\subsection{Definición y Fundamentos}

En el CAPM, dado que todos los agentes poseen carteras bien diversificadas, van a exigir una prima en función del riesgo sistemático de cada activo, y no del riesgo específico. Como el riesgo sistemático se mide por la beta, entonces, la rentabilidad exigida será función de la beta.

\begin{definicion}[Línea del Mercado de Títulos (SML)]
A la recta que representa la relación entre rentabilidad esperada de todos los activos en función de su riesgo sistemático se le conoce como \textbf{Línea del Mercado de Activos} (SML, Security Market Line). La SML muestra la rentabilidad exigida (en equilibrio) a cada activo como función de su riesgo sistemático, medido por su beta.
\end{definicion}

La ecuación fundamental del CAPM nos dice que la prima de riesgo de un activo individual será función de:
\begin{itemize}
\item La prima de riesgo esperada del mercado
\item El riesgo sistemático de ese activo (su beta)
\end{itemize}

\subsection{La Ecuación del CAPM}

\begin{teorema}[Ecuación Fundamental del CAPM]
La rentabilidad esperada de un activo $i$ según el CAPM viene dada por:

\begin{equation}
E[r_i] = r_f + \beta_i(E[r_M] - r_f)
\end{equation}

donde:
\begin{itemize}
\item $E[r_i]$ es la rentabilidad esperada del activo $i$
\item $r_f$ es la tasa libre de riesgo
\item $\beta_i$ es la beta del activo $i$
\item $E[r_M]$ es la rentabilidad esperada del mercado
\item $(E[r_M] - r_f)$ es la prima de riesgo del mercado
\end{itemize}
\end{teorema}

En este tema vamos a representar en mayúsculas las primas de riesgo, así:
\begin{align}
r_M - r_f &= R_M \rightarrow \text{Prima de riesgo del mercado} \\
r_i - r_f &= R_i \rightarrow \text{Prima de riesgo del activo } i
\end{align}

Por tanto, la ecuación del CAPM también puede expresarse como:

\begin{equation}
E[R_i] = \beta_i E[R_M]
\end{equation}

donde $E[R_i]$ es la prima de riesgo esperada del activo $i$ y $E[R_M]$ es la prima de riesgo esperada del mercado.

\begin{ejemplo}[Aplicación del CAPM]
Suponga que se cumplen todos los supuestos del modelo CAPM, y deseamos conocer la rentabilidad esperada de las acciones de AMADEUS sabiendo que:
\begin{itemize}
\item La rentabilidad esperada del mercado para el próximo año es del 11.5\%
\item La rentabilidad ofrecida por las letras del tesoro a un año es del 3.5\%
\item La beta de las acciones de AMADEUS es de 1.8
\end{itemize}

\textbf{Pregunta 1:} ¿Será la rentabilidad esperada mayor que la del mercado?

Dado que $\beta_i = 1.8 > 1$, la rentabilidad esperada de AMADEUS será mayor que la del mercado, ya que el activo tiene un riesgo sistemático superior al del mercado.

\textbf{Pregunta 2:} Determine la rentabilidad esperada de las acciones de AMADEUS.

Aplicando la ecuación del CAPM:
\begin{align}
E[r_i] &= r_f + \beta_i(E[r_M] - r_f) \\
E[r_i] &= 0.035 + 1.8(0.115 - 0.035) \\
E[r_i] &= 0.035 + 1.8 \times 0.08 \\
E[r_i] &= 0.035 + 0.144 \\
E[r_i] &= 0.179 = 17.9\%
\end{align}

La rentabilidad esperada de las acciones de AMADEUS es del 17.9\%.
\end{ejemplo}

La pendiente de la SML es igual a la prima de riesgo del mercado $(E[r_M] - r_f)$. Un activo con beta igual a 1 tendrá una rentabilidad esperada igual a la del mercado. Un activo con beta mayor que 1 tendrá una rentabilidad esperada mayor que la del mercado, y un activo con beta menor que 1 tendrá una rentabilidad esperada menor que la del mercado.

\textbf{Pregunta:} ¿Puede algún activo situarse fuera de la SML?

En equilibrio, todos los activos deben situarse sobre la SML. Si un activo se situara por encima de la SML, ofrecería una rentabilidad esperada mayor que la que corresponde a su nivel de riesgo sistemático, lo que generaría una demanda excesiva que elevaría su precio y reduciría su rentabilidad esperada hasta situarlo sobre la SML. De manera análoga, si un activo se situara por debajo de la SML, ofrecería una rentabilidad esperada menor que la que corresponde a su nivel de riesgo sistemático, lo que generaría una oferta excesiva que reduciría su precio y elevaría su rentabilidad esperada hasta situarlo sobre la SML. Por tanto, en equilibrio de mercado, todos los activos deben situarse exactamente sobre la SML.

\section{La Beta}

\subsection{Definición y Significado}

En el tema anterior hemos estudiado que el riesgo total de un activo puede descomponerse en riesgo específico (diversificable) y riesgo sistemático (no diversificable). La beta mide el riesgo sistemático de un activo.

\begin{definicion}[Beta]
La \textbf{beta} ($\beta$) de un activo mide la contribución de ese activo al riesgo de una cartera bien diversificada o a la cartera de mercado. La beta indica la sensibilidad de la rentabilidad en exceso de un activo individual \(i\) ante movimientos de la rentabilidad del mercado.
\end{definicion}

La beta se puede calcular de dos formas equivalentes:

\begin{teorema}[Cálculo de la Beta]
La beta de un activo $i$ puede calcularse mediante:

\begin{equation}
\beta_i = \frac{\sigma_{i,M}}{\sigma_M^2} = \frac{\text{Cov}(R_i, R_M)}{\text{Var}(R_M)}
\end{equation}

donde:
\begin{itemize}
\item $\sigma_{i,M}$ es la covarianza entre los rendimientos del activo $i$ y del mercado
\item $\sigma_M^2$ es la varianza de los rendimientos del mercado
\item $\text{Cov}(R_i, R_M)$ es la covarianza entre las primas de riesgo del activo $i$ y del mercado
\item $\text{Var}(R_M)$ es la varianza de la prima de riesgo del mercado
\end{itemize}
\end{teorema}

\subsection{Cálculo de Beta mediante Regresión}

La beta también se puede calcular a través de una regresión entre la rentabilidad en exceso del activo y del mercado.

\begin{proposicion}[Regresión de Mercado]
La beta puede estimarse mediante la regresión:

\begin{equation}
R_{i,t} = \alpha_i + \beta_i R_{M,t} + \varepsilon_{i,t}
\end{equation}

donde:
\begin{itemize}
\item $R_{i,t} = r_{i,t} - r_f$ es la rentabilidad en exceso del activo $i$ en el periodo $t$
\item $R_{M,t} = r_{M,t} - r_f$ es la rentabilidad en exceso del mercado en el periodo $t$
\item $\alpha_i$ es el intercepto (alfa de Jensen)
\item $\beta_i$ es la pendiente (beta del activo)
\item $\varepsilon_{i,t}$ es el término de error
\end{itemize}
\end{proposicion}

El coeficiente de pendiente $\beta_i$ de esta regresión es precisamente la beta del activo, ya que:

\begin{equation}
\beta_i = \frac{\text{Cov}(R_i, R_M)}{\text{Var}(R_M)}
\end{equation}

\begin{ejemplo}[Cálculo de Beta]
Determine la beta y la rentabilidad esperada de las acciones de la empresa TELEFON sabiendo que:
\begin{itemize}
\item La covarianza entre los rendimientos de la empresa y del IBEX es de 0.0099
\item La desviación típica de la empresa es del 17\%
\item La desviación típica del índice de mercado (IBEX) es del 24\%
\item La rentabilidad de las letras del tesoro es del 4.5\%
\item La prima de riesgo esperada del mercado es del 8\%
\end{itemize}

\textbf{Solución:}

Primero, calculamos la beta:
\begin{align}
\beta_i &= \frac{\text{Cov}(R_i, R_M)}{\text{Var}(R_M)} \\
\beta_i &= \frac{0.0099}{(0.24)^2} \\
\beta_i &= \frac{0.0099}{0.0576} \\
\beta_i &= 0.1719
\end{align}

Ahora, calculamos la rentabilidad esperada usando el CAPM:
\begin{align}
E[r_i] &= r_f + \beta_i(E[r_M] - r_f) \\
E[r_i] &= 0.045 + 0.1719 \times 0.08 \\
E[r_i] &= 0.045 + 0.01375 \\
E[r_i] &= 0.05875 = 5.875\%
\end{align}

La beta de TELEFON es 0.1719 y su rentabilidad esperada es del 5.875\%.
\end{ejemplo}

\subsection{Clasificación de Acciones según su Beta}

Podemos distinguir entre diferentes tipos de acciones por su beta:

\begin{definicion}[Acciones Defensivas]
Las \textbf{acciones defensivas} son aquellas con $\beta < 1$. Estas acciones tienen un riesgo sistemático menor que el del mercado y, por tanto, su rentabilidad esperada es menor que la del mercado. Son menos sensibles a los movimientos del mercado.
\end{definicion}

\begin{definicion}[Acciones Neutras]
Las \textbf{acciones neutras} son aquellas con $\beta = 1$. Estas acciones tienen un riesgo sistemático igual al del mercado y, por tanto, su rentabilidad esperada es igual a la del mercado.
\end{definicion}

\begin{definicion}[Acciones Agresivas]
Las \textbf{acciones agresivas} son aquellas con $\beta > 1$. Estas acciones tienen un riesgo sistemático mayor que el del mercado y, por tanto, su rentabilidad esperada es mayor que la del mercado. Son más sensibles a los movimientos del mercado.
\end{definicion}

\begin{proposicion}[Beta del Mercado y del Activo Libre de Riesgo]
\begin{itemize}
\item La beta del mercado es igual a 1: $\beta_M = 1$
\item La beta del activo libre de riesgo es igual a 0: $\beta_{rf} = 0$
\end{itemize}
\end{proposicion}

\begin{proof}
Para el mercado:
\begin{equation}
\beta_M = \frac{\text{Cov}(R_M, R_M)}{\text{Var}(R_M)} = \frac{\text{Var}(R_M)}{\text{Var}(R_M)} = 1
\end{equation}

Para el activo libre de riesgo, como su rentabilidad es constante e igual a $r_f$, su covarianza con el mercado es cero:
\begin{equation}
\beta_{rf} = \frac{\text{Cov}(r_f, R_M)}{\text{Var}(R_M)} = \frac{0}{\text{Var}(R_M)} = 0
\end{equation}
\end{proof}

\section{La Beta de una Cartera}

\subsection{Cálculo de la Beta de una Cartera}

Dado que la beta mide el riesgo no diversificable, la beta de una cartera es simplemente la suma de las betas de cada activo ponderadas por el peso de cada activo en la cartera.

\begin{teorema}[Beta de una Cartera]
La beta de una cartera es la media ponderada de las betas de los activos que la componen:

\begin{equation}
\beta_p = \sum_{i=1}^{N} w_i \beta_i
\end{equation}

donde:
\begin{itemize}
\item $\beta_p$ es la beta de la cartera
\item $w_i$ es el peso del activo $i$ en la cartera
\item $\beta_i$ es la beta del activo $i$
\item $N$ es el número de activos en la cartera
\end{itemize}
\end{teorema}

\begin{proof}
La rentabilidad de la cartera es:
\begin{equation}
r_p = \sum_{i=1}^{N} w_i r_i
\end{equation}

La prima de riesgo de la cartera es:
\begin{equation}
R_p = r_p - r_f = \sum_{i=1}^{N} w_i (r_i - r_f) = \sum_{i=1}^{N} w_i R_i
\end{equation}

Aplicando la definición de beta:
\begin{align}
\beta_p &= \frac{\text{Cov}(R_p, R_M)}{\text{Var}(R_M)} \\
&= \frac{\text{Cov}\left(\sum_{i=1}^{N} w_i R_i, R_M\right)}{\text{Var}(R_M)} \\
&= \frac{\sum_{i=1}^{N} w_i \text{Cov}(R_i, R_M)}{\text{Var}(R_M)} \\
&= \sum_{i=1}^{N} w_i \frac{\text{Cov}(R_i, R_M)}{\text{Var}(R_M)} \\
&= \sum_{i=1}^{N} w_i \beta_i
\end{align}
\end{proof}

Para el caso particular de dos activos:

\begin{equation}
\beta_p = w_1 \beta_1 + w_2 \beta_2
\end{equation}

\begin{ejemplo}[Cálculo de la Beta de una Cartera]
Un fondo de inversión tiene una cartera compuesta por tres activos con riesgo con las siguientes características:
\begin{itemize}
\item Activo 1: $\beta_1 = 0.05$, peso $w_1 = 20\%$
\item Activo 2: $\beta_2 = 1.02$, peso $w_2 = 35\%$
\item Activo 3: Covarianza con el mercado = 0.0399, desviación típica del mercado = 19\%, peso $w_3 = 45\%$
\end{itemize}

Calcule la beta de la cartera del fondo de inversión.

\textbf{Solución:}

Primero, calculamos la beta del activo 3:
\begin{align}
\beta_3 &= \frac{\text{Cov}(R_3, R_M)}{\text{Var}(R_M)} \\
\beta_3 &= \frac{0.0399}{(0.19)^2} \\
\beta_3 &= \frac{0.0399}{0.0361} \\
\beta_3 &= 1.1058
\end{align}

Ahora, calculamos la beta de la cartera:
\begin{align}
\beta_p &= w_1 \beta_1 + w_2 \beta_2 + w_3 \beta_3 \\
\beta_p &= 0.20 \times 0.05 + 0.35 \times 1.02 + 0.45 \times 1.1058 \\
\beta_p &= 0.01 + 0.357 + 0.4976 \\
\beta_p &= 0.8646
\end{align}

La beta de la cartera es 0.8646.
\end{ejemplo}

\section{Extensiones y Otros Modelos de Valoración}

\subsection{Estudios de Eventos (Event Studies)}

Los estudios de eventos constituyen una metodología empírica ampliamente utilizada en finanzas para analizar el impacto de eventos corporativos específicos sobre el valor para el accionista.

\begin{definicion}[Estudio de Eventos]
Un \textbf{estudio de eventos} (event study) es una metodología que analiza la reacción del mercado ante el anuncio de decisiones empresariales o eventos específicos, midiendo los rendimientos anormales alrededor de la fecha del evento para evaluar el impacto sobre el valor para el accionista.
\end{definicion}

Ejemplos de eventos que pueden estudiarse incluyen:
\begin{itemize}
\item Dividendos
\item Emisiones de acciones
\item Emisiones de deuda
\item Splits (desdoblamientos) de acciones
\item Fusiones y adquisiciones
\end{itemize}

Si las empresas en cuestión cotizan en bolsa, se puede observar la reacción del mercado al anuncio de que se tomará una determinada decisión. Si el mercado entiende las consecuencias de esa decisión para el valor del accionista, entonces la reacción del precio de la acción reflejará el valor añadido para el accionista.

\subsubsection{Metodología de los Estudios de Eventos}

Los estudios de eventos se consideran una forma de poner a prueba la eficiencia del mercado: si los mercados financieros son eficientes en información, debería producirse una reacción inmediata al evento en la fecha del anuncio, y no debería haber más reacción en los días posteriores de cotización.

Es habitual analizar un periodo alrededor de la fecha del evento: una ventana de corto plazo, que abarca aproximadamente diez días de cotización antes del evento y hasta diez días después. MacKinlay (1997) muestra que la capacidad de los estudios de eventos para detectar rendimientos anormales es mucho mayor cuando se emplean datos diarios, en lugar de datos mensuales, trimestrales o anuales. Sin embargo, los datos intradiarios suelen estar llenos de ruido de microestructura.

\subsubsection{Modelo de Mercado}

Para estimar el rendimiento normal, podemos utilizar el modelo de mercado:

\begin{equation}
r_{i,t} = \alpha_i + \beta_i r_{M,t} + \varepsilon_{i,t}
\end{equation}

donde $r_{M,t}$ es el rendimiento del mercado (por ejemplo, de un índice) y $r_{i,t}$ es el rendimiento de la empresa.

Para cada empresa de la muestra, se recopilan datos de rendimientos durante, por ejemplo, 200 días antes del evento (es decir, antes del anuncio). Luego se estiman los parámetros $\alpha_i$ y $\beta_i$ mediante regresión OLS (mínimos cuadrados ordinarios). Armitage (1995) sugiere que los periodos de estimación pueden abarcar entre 100 y 300 días cuando se usan datos diarios, y entre 24 y 60 meses cuando el análisis se realiza con datos mensuales.

El rendimiento esperado normal para la empresa $i$ en la fecha del anuncio ($t = 0$) se calcula como:

\begin{equation}
E[r_{i,0}] = \hat{\alpha}_i + \hat{\beta}_i r_{M,0}
\end{equation}

El rendimiento anormal estimado es:

\begin{equation}
AR_i = r_{i,0} - E[r_{i,0}] = r_{i,0} - (\hat{\alpha}_i + \hat{\beta}_i r_{M,0})
\end{equation}

El rendimiento anormal es la diferencia entre el rendimiento real de la acción y el rendimiento esperado según el modelo de mercado.

Una vez que se ha recopilado el rendimiento anormal de cada empresa de la muestra, se pueden responder varias preguntas interesantes:
\begin{itemize}
\item ¿Cuál es el rendimiento anormal promedio y es significativamente distinto de cero? (Se calcula el promedio y se realiza una prueba $t$ para contrastar su significancia).
\item ¿Existe heterogeneidad en las respuestas de las empresas? ¿Qué características de las firmas determinan esas diferencias en las reacciones?
\item ¿Cuáles son los determinantes de los rendimientos del anuncio? (Se pueden regresar los rendimientos anormales sobre variables específicas de la empresa u otras variables explicativas).
\end{itemize}

\subsubsection{Aplicaciones de los Estudios de Eventos}

Los estudios de eventos se han aplicado ampliamente en la práctica financiera. Por ejemplo, NBIM (Norges Bank Investment Management) utilizó esta metodología en 2012 para analizar el impacto de las inversiones de fondos soberanos en los mercados financieros. Los datos de factores de riesgo utilizados en estos estudios, como los factores de Fama-French y Carhart, están disponibles públicamente en la biblioteca de datos de Kenneth French: \url{http://mba.tuck.dartmouth.edu/pages/faculty/ken.french/data_library.html}.

\subsection{Modelo de Tres Factores de Fama-French}

El modelo de tres factores de Fama y French fue desarrollado por Eugene Fama y Kenneth French en 1993 para describir los rendimientos de las acciones. El CAPM utiliza solo una variable para explicar los rendimientos de una cartera o acción: el rendimiento del mercado. El modelo de Fama y French, en cambio, emplea tres variables.

\begin{teorema}[Modelo de Tres Factores de Fama-French]
Fama y French observaron que dos tipos de acciones tienden a obtener rendimientos superiores al promedio del mercado:
\begin{enumerate}
\item Las empresas de pequeña capitalización (small caps)
\item Las acciones con una alta relación valor contable / valor de mercado (book-to-market ratio)
\end{enumerate}

Por ello, añadieron dos factores adicionales al CAPM para reflejar la exposición de una cartera a estas dos clases de acciones.
\end{teorema}

El modelo de tres factores de Fama-French se expresa como:

\begin{equation}
r_{i,t} = \alpha_i + \beta_1 r_{m,t} + \beta_2 \text{SMB}_t + \beta_3 \text{HML}_t + \varepsilon_{i,t}
\end{equation}

donde:
\begin{itemize}
\item $r_{i,t}$ es el rendimiento del activo $i$ en el periodo $t$
\item $\alpha_i$ es el intercepto (alfa de Jensen)
\item $\beta_1$ es la sensibilidad al factor de mercado
\item $r_{m,t}$ es el rendimiento del mercado en el periodo $t$
\item $\beta_2$ es la sensibilidad al factor de tamaño (SMB)
\item $\text{SMB}_t$ (Small Minus Big) es la prima mensual del factor tamaño
\item $\beta_3$ es la sensibilidad al factor valor (HML)
\item $\text{HML}_t$ (High Minus Low) es la prima mensual del factor valor contable/mercado
\item $\varepsilon_{i,t}$ es el término de error
\end{itemize}

\begin{definicion}[Factor SMB]
El \textbf{SMB} (Small Minus Big) es una cartera de inversión de costo cero que mantiene posiciones largas en acciones de pequeña capitalización y cortas en acciones de gran capitalización.
\end{definicion}

\begin{definicion}[Factor HML]
El \textbf{HML} (High Minus Low) es una cartera de inversión de costo cero que mantiene posiciones largas en acciones con alto valor contable/mercado y cortas en las de bajo valor contable/mercado.
\end{definicion}

\subsection{Modelo de Cuatro Factores de Carhart}

El modelo de cuatro factores de Carhart (1997) es una extensión del modelo de Fama-French que incluye un factor adicional de momentum.

\begin{teorema}[Modelo de Cuatro Factores de Carhart]
El modelo de Carhart añade un cuarto factor de momentum al modelo de tres factores de Fama-French:

\begin{equation}
r_{i,t} = \alpha_i + \beta_1 r_{m,t} + \beta_2 \text{SMB}_t + \beta_3 \text{HML}_t + \beta_4 \text{UMD}_t + \varepsilon_{i,t}
\end{equation}

donde $\text{UMD}_t$ (Up Minus Down) es la prima mensual del factor momentum.
\end{teorema}

\begin{definicion}[Factor UMD]
El \textbf{UMD} (Up Minus Down) es una cartera de costo cero, larga en las acciones ganadoras de los últimos 12 meses y corta en las acciones perdedoras, según la definición del sitio web de Kenneth French.
\end{definicion}

El modelo regresa los rendimientos en exceso del activo sobre un intercepto (alfa) y varios factores en el lado derecho de la ecuación que intentan controlar los riesgos de mercado generales:
\begin{enumerate}
\item Rendimiento mensual del índice de mercado menos la tasa libre de riesgo
\item Prima mensual del factor: valor contable/valor de mercado (HML)
\item Prima mensual del factor: tamaño (SMB)
\item Prima mensual del factor: momentum (UMD), es decir, ganadores menos perdedores
\end{enumerate}

Un inversor demuestra capacidad predictiva o habilidad de gestión cuando su cartera presenta un alfa positivo y estadísticamente significativo.

\subsection{Teoría de Valoración por Arbitraje (APT)}

La Teoría de Valoración por Arbitraje (APT) fue propuesta por Stephen Ross en 1976 y constituye un modelo alternativo de valoración de activos.

\begin{definicion}[Teoría de Valoración por Arbitraje (APT)]
La \textbf{Teoría de Valoración por Arbitraje} (APT, Arbitrage Pricing Theory) es un modelo de valoración de activos que parte de la idea de que si dos activos tienen los mismos riesgos sistemáticos, deben ofrecer el mismo rendimiento esperado. Si no es así, habría una oportunidad de arbitraje (ganancia sin riesgo), y los precios se ajustarían hasta eliminarla.
\end{definicion}

El APT dice que el rendimiento de un activo puede explicarse por varios factores de riesgo sistemático, no solo por el mercado como en el CAPM.

\begin{teorema}[Ecuación del APT]
Según el APT, el rendimiento esperado de un activo viene dado por:

\begin{equation}
E[R_i] = R_f + b_{i1}\lambda_1 + b_{i2}\lambda_2 + \cdots + b_{ik}\lambda_k
\end{equation}

donde:
\begin{itemize}
\item $E[R_i]$ es la rentabilidad esperada del activo $i$
\item $R_f$ es la tasa libre de riesgo
\item $b_{ij}$ es la exposición (factor loading) del activo $i$ al factor de riesgo $j$
\item $\lambda_j$ es la prima de riesgo asociada al factor $j$
\item $k$ es el número de factores de riesgo sistemático
\end{itemize}
\end{teorema}

Cada factor ($\lambda_j$) representa una fuente de riesgo sistemático (por ejemplo, crecimiento del PIB, inflación, cambios en tasas de interés, primas de liquidez, etc.). Cada activo tiene una exposición distinta ($b_{ik}$) a esos factores. Los inversores son compensados por asumir esos riesgos, pero no por riesgos diversificables.

El APT amplía la visión del CAPM al reconocer que los rendimientos dependen de múltiples riesgos sistemáticos. El CAPM supone un único factor (el mercado), mientras que el APT permite varios factores y se basa solo en la idea de que no puede haber oportunidades de arbitraje persistentes.

\section{Conclusiones}

Este documento ha presentado un tratamiento riguroso de los modelos de valoración de activos con riesgo, con especial énfasis en el CAPM. El CAPM constituye un marco teórico fundamental que permite entender cómo se determinan los precios de los activos financieros en equilibrio de mercado, estableciendo una relación directa entre la rentabilidad esperada de un activo y su riesgo sistemático, medido por la beta.

La línea del mercado de capitales (CML) muestra las combinaciones eficientes entre el activo libre de riesgo y la cartera de mercado, mientras que la línea del mercado de títulos (SML) establece la relación entre la rentabilidad esperada y el riesgo sistemático de cada activo individual.

El concepto de beta como medida del riesgo sistemático es central en el modelo CAPM y permite clasificar los activos según su sensibilidad a los movimientos del mercado. La beta de una cartera es simplemente la media ponderada de las betas de los activos que la componen.

Las extensiones del CAPM, como el modelo de tres factores de Fama-French, el modelo de cuatro factores de Carhart y la Teoría de Valoración por Arbitraje (APT), reconocen que los rendimientos pueden depender de múltiples factores de riesgo sistemático, proporcionando modelos más ricos y realistas para la valoración de activos.

Los estudios de eventos constituyen una metodología empírica importante que permite evaluar el impacto de decisiones corporativas sobre el valor para el accionista y poner a prueba la eficiencia informacional de los mercados financieros.

A pesar de sus limitaciones y supuestos restrictivos, el CAPM y sus extensiones continúan siendo herramientas fundamentales en la práctica financiera moderna, proporcionando un marco de referencia esencial para la valoración de activos, la gestión de carteras y la evaluación del rendimiento de inversiones.

\end{document}

